\section{Introduction}\label{s:intro}
Data centres consumed approximately 415 TWH, or 1.5\% of global energy consumption in 2024, and this is expected to more than double by 2030 \cite{iea_energy_2025}.
Cloud computing happens largely in data centers and significantly contributes to this number, which does not even take the water use \cite{siddik_environmental_2021} into account. The same goes for other far-reaching environmental impacts that remain obscured due to the largely decentralized and commercial nature of the cloud \cite{Monserrate2022Cloud}. According to various surveys, more than 70\% of IT organisations use microservices \cite{statista_global_2021, gartner_peer_community_microservices_2023, loukides_microservices_2020}.
This suggests that not just cloud computing in general, but microservice applications in particular, form relevant objects of study with regards to their energy usage \cite{araujo_energy_2024}. While microservices have been widely adopted because the loosely coupled, small and lightweight components are scalable, flexible and can be used to increase resilience, their deployment also has caveats, such as increased overhead and complexity \cite{pirhonen_pains_2025}. In order to understand and reduce the resource utilization and environmental impact of microservices, realistic simulation of such architectures in cloud datacenters is important \cite{katal_energy_2023}. Through simulation, potential areas of improvement to reduce energy usage can be identified and strategies to reduce resource consumption can be tried out at lower cost \cite{mansouri_cloud_2020}. A large number of cloud simulators is available, but they all differ in the capabilities they offer, components they model, their implementation and ease of use.

In this study, we compare five cloud simulators, \textit{CloudSim},\textit{DISSECT-CF}, \textit{NetworkCloudSim}, \textit{OpenDC} and \textit{SimGrid} that offer the capability to simulate energy usage\footnote{The replication package for this study can be found at \url{https://github.com/S2-group/simpat-2026-cloud-simulators-energy}.}. These simulators were selected by extracting 67 simulators from 23 peer-reviewed articles reviewing cloud simulators and choosing the five most recently updated simulators that have energy usage simulation capabilities and for which the code is publicly accessible. 
\rone{Notably, none of the considered simulators provides support for modeling energy consumption at the level of individual microservices. As a result, the scope of this study is limited to the estimation and evaluation of total host-level energy consumption.}
Then, we assess the simulators according to the accuracy of the simulated \rone{host-level} energy values they produce for traces of two different microservice applications, \textit{Sock Shop}\footnote{\url{https://github.com/microservices-demo/microservices-demo}} and \textit{Train Ticket} \cite{zhou_benchmarking_2018} by comparing the results to real-world measures. Furthermore, we describe the extent to which the simulators support energy-related simulation through what resources can be simulated but also what other capabilities related to energy are offered. While the software systems the simulators are set to simulate are microservices, we also look at the affordances of the simulators themselves and how well they are suited to simulating microservice applications.

Our results show that \textit{OpenDC} is the simulator most suited to simulating \rone{host-level} microservice applications, both in terms of accuracy of its energy usage simulations and its suitability to simulating microservice architecture style applications. It has both the lowest minimum (0.14) and maximum (0.41) mean relative error and is the second-fastest simulator with simulation times of less than one second. \textit{NetworkCloudSim} has both the highest minimum (0.42) and maximum (0.95) mean relative error. \textit{SimGrid} without simulating communication is the fastest simulator with a mean simulation time of 0.04 seconds, while it is the slowest (12.02 seconds on average) when including communication in its simulation. Depending on the context such as workload definition, user skill and preference, and desired levels of configurability, other simulators might be well suited as well. By doing so, we enable prospective users of cloud simulators the possibility to make an informed choice about what cloud simulator to use, and perhaps even make it more attractive to start to do so. Furthermore, making a more informed choice in cloud simulator could lead to more accurate simulation results as we show that the simulator currently most used, also in the context of microservices \cite{araujo_energy_2024}, \textit{CloudSim}, is not always the best choice to make. Researchers and engineers working with microservice systems can benefit from this study as it can help them with making the choice of cloud simulator easier and more informed, in turn enabling them to design and develop more energy-conscious microservice systems.


